\chapter{Why local evaluation can be exact}\label{ch:local}
The project asks for local valuation without a global observer deciding the
worth of every action. The mathematical possibility comes from cancellation,
not from approximating a large world by a small one. A local calculation is
exact when it includes every potential term that can change.

\section{The separable theorem}
Let $V(x)=\sum_i v_i(x_i)$ and let $W$ contain all coordinates changed by
$\delta$. For $i\notin W$, $z_i+\delta_i=z_i$. Split both sums:
\begin{align}
 V(z)-V(z+\delta)
 &=\sum_{i\in W}[v_i(z_i)-v_i(z_i+\delta_i)]\\
 &\quad+\sum_{i\notin W}[v_i(z_i)-v_i(z_i)]\\
 &=V_W(z)-V_W(z+\delta).
\end{align}
No information about an unchanged term is needed to calculate its zero
difference. This identity holds for nonlinear $v_i$ and does not require
Gaussian geometry. The Gaussian model is one especially transparent
specialization of the historical locality result \cite{foundation}.

\section{A complete small example}
Use references $(10,10,10)$ and scales $(2,2,2)$ at state $(18,2,12)$.
The first two potential terms are eight each; the third is one-half.
A four-unit transfer from cell one to cell two gives $(14,6,12)$, with
potential terms two, two and one-half. The global difference is
$16.5-4.5=12$. The two-endpoint local calculation is $16-4=12$.
The remote cell's value was not assumed zero. It cancelled because it
was unchanged.

If an action also changes cell three, the proof demands its inclusion.
Calling the third effect small or inconvenient does not permit its omission
while retaining the adjective exact. One may define an approximation, but
then its error and claim must be stated separately.

\section{The factor version}
A more general potential has factors
\begin{equation}
 V(x)=\sum_{\alpha\in\mathcal F}\phi_\alpha(x_{A_\alpha}),
\end{equation}
where $A_\alpha$ is the set of arguments used by factor $\alpha$. Let
$\mathcal F_W=\{\alpha:A_\alpha\cap W\ne\varnothing\}$.
Every factor outside $\mathcal F_W$ cancels, so the exact difference is
the sum over $\mathcal F_W$. A factor in that set may require reading a
coordinate outside $W$. Write support and read support need not coincide.

For example a declared term $k(x_1-x_3)^2$ changes when an action changes
cell one. Cell three may remain physically untouched, yet its current
value is required to evaluate the factor. A two-endpoint view that omits
it cannot be called complete for that expanded potential. This is an
illustrative extension, not an adopted nonseparable Gaussian model.

\section{Three different locality claims}
\emph{Valuation locality} means the finite value needs only complete
affected factors. \emph{Information locality} means a particular actor or
module may access only specified records. \emph{Coordination locality}
concerns how candidate groups are assembled and selected. A proof of the
first does not establish the second or third.

A central program could enumerate every group in the world and evaluate
each group locally. Its values would satisfy cancellation, but the selector
would still be global. Similarly a class named LocalView is not a complete
security proof if a helper can read unrestricted state elsewhere. Part III
turns these distinctions into review questions without asserting that the
future runtime already enforces them.

\section{Frozen state and version completeness}
Locality also requires agreement about time and parameters. Combining the
source's present stock with the destination's stale stock does not evaluate
one declared state. Changing a scale between the before and after terms
does not compute one potential difference. Exact arithmetic cannot repair
either mismatch.

Thus a reconstructible local quote binds the state version, affected factor
identities, reference/scale version, action increment and process-burden
declaration. Hashes can identify those inputs. They cannot certify that a
sensor observed reality correctly; that is a measurement question.

\section*{Exercise and answer}
Suppose a quote reports a difference of twelve while a global audit reports
eleven. Name three possible causes before blaming numerical precision.
An affected factor may be missing; the states may belong to different
epochs; or one calculation may use a different parameter version. Each is
a semantic discrepancy. The exact identity tells the reviewer where to
look because, under its actual assumptions, such a difference is impossible.

\section*{Chapter recap}
Local exactness is a cancellation theorem over complete affected factors.
It is powerful precisely because its boundary is precise. It does not
certify sensors, ownership, or decentralized selection.
